\chapter{Exporter une image}
\label{chap:export}
\section{Choisir un format}
Ouvrir la famille Export et choisir le format. Le contrôle Qualité JPEG
n’apparaît que pour JPEG ; il va de 1 à 100, référence 92.
\begin{longtable}{@{}L{2cm}L{4.6cm}L{7cm}@{}}
\toprule Format & Encodage actuel & Usage et limites\\\midrule\endhead
JPEG & sRGB, 8 bits & Copie courante ; compression avec perte, qualité réglable.\\
PNG & sRGB, 16 bits & Image rendue à haute profondeur ; pas de réglage Qualité JPEG.\\
TIFF & sRGB, 16 bits & Référence de comparaison ou échange de rendu à haute profondeur.\\
DNG & RGB linéaire, 16 bits, primaires sRGB & Contient le rendu développé ; ne préserve pas la mosaïque RAW d’origine.\\
\bottomrule
\end{longtable}
Ces propriétés décrivent l’exporteur actuel, pas toutes les capacités théoriques
de chaque format. Aucun sélecteur utilisateur de gamut de sortie ou de redimensionnement
n’est proposé dans ce panneau.

\section{Destination et confirmation}
\begin{enumerate}
\item Attendre la fin du développement et cliquer sur \emph{Choisir la destination…}.
\item Vérifier le nom, le dossier et le type du dialogue natif. Le format du
panneau est sélectionné initialement ; changer le type dans le dialogue change
aussi le format demandé.
\item Dans \emph{Vérifier l’export}, contrôler format, encodage, dimensions
pleine résolution et destination demandée.
\item Cliquer sur \emph{Exporter}, ou \emph{Annuler} pour revenir sans écrire.
\item Lire le message final pour connaître le fichier réellement publié ou l’erreur.
\end{enumerate}
L’extension est adaptée au format. Si le nom est pris, l’exporteur ajoute
un numéro disponible plutôt que d’écraser le fichier. Le chemin final peut
donc différer du nom demandé dans le récapitulatif.

\attention{Le récapitulatif ne promet pas un nom final réservé. Un autre
processus peut créer un fichier entre le choix et l’écriture ; la protection
contre les collisions reste assurée par l’exporteur.}

\section{Gestes et contrôles}
\touche{Tab} permet de parcourir les commandes du dialogue et \touche{Entrée}
de valider une saisie active. Utiliser explicitement le bouton Exporter pour
lancer l’écriture ; aucun \touche{Ctrl + E} global n’est déclaré dans IPR.
Ne pas confondre l’export avec la sauvegarde automatique des compagnons.

Pour un test reproductible, exporter une référence neutre et l’état retouché
avec le même profil, la même balance, les mêmes corrections optiques et le
même format. Noter les réglages : un TIFF IPR ne constitue pas à lui seul
un relevé complet et fiable de la recette appliquée.
